\section{Описание}
Нужно разработать программу для токенизации текстовых документов, которая, разобьет текст на токены, приведет их к нижнему регистру и удалит стоп-слова. Токенизация это первый этап в процессе индексации документов для информационного поиска. Нужно разработать правила разбиения текста на токены, реализовать алгоритм токенизации, провести анализ результатов работы и собрать статистику: количество токенов, среднюю длину токена, время выполнения программы и скорость обработки данных.

\pagebreak

\section{Исходный код}
Для разработки программ использовался язык программирования C. В конечном итоге был получен набор утилит:

\begin{enumerate}
    \item \textbf{./tokenizer <in\_dir> <out\_dir>} –- разбивает текст на отдельные токены
    \item \textbf{./normalizer <in\_dir> <out\_dir>} -- приводит текст к нижнему регистру и заменяет букву «ё» на «е»
    \item \textbf{./stopwords <in\_dir> <out\_dir>} -- удаляет стоп-слова (неинформативные слова)
    \item \textbf{./pipeline <in\_dir> <out\_dir>} -- применяет весь пайплайн обработки одновременно, а также собирает статистику
\end{enumerate}

Были написаны небольшие библиотеки для токенизации, нормализации и удаления стоп слов из текста. В процессе отладки было удобно применять эти этапы поочередно, чтобы верифицировать результат, поэтому были выделены отдельные утилиты для каждого этапа, а также общая, применяющая весь пайплайн обработки к текстам сразу и собирающая статистику.

\vspace{1 em}
\noindent\textbf{Работа с UTF-8.} UTF‑8 -- это многобайтовая кодировка, где один символ может занимать от 1 до 4 байтов. Если читать файл как обычную строку char*, мы не можем напрямую обращаться к символам по индексу (например, text[i]), потому что один символ может быть разбит на несколько байтов. \\
Поэтому был написан набор функций для преобразования входного набора байт в массив wchar\_t. Тип wchar\_t представляет собой широкий символ, который хранит каждый символ Юникода отдельно. Это позволило корректно обращаться к символам по индексу, а также использовать функции для определения длины текста, нахождения символов (wcslen, wcschr).

\vspace{1 em}
\noindent\textbf{Токенизация.} Библиотека \texttt{tokenizer\_lib.c} реализует функцию токенизации текста на широких символах wchar\_t. Токенизация выполняется путем разделения текста на токены с удалением разделителей и заменой их одиночными пробелами между токенами.

\begin{enumerate}
    \item \textbf{Определяем разделители:} пробелы, знаки пунктуации и специальные символы
    
    \item \textbf{Проходим по тексту} символ за символом
    
    \item \textbf{Находим токены:}
    \begin{itemize}
        \item Все между разделителями = один токен
        \item Последовательность разделителей = один пробел между токенами
        \item Разделители на краях отбрасываются
    \end{itemize}
    
    \item \textbf{Сохраняем результат:} токены через один пробел, без лишних пробелов
\end{enumerate}

\begin{lstlisting}[language=C]
static
bool is_delimiter( wchar_t c )
{
    if ( iswspace( c ) || iswpunct( c ) )
        return 1;

    switch ( c )
    {
        case L'«': case L'»': case L'„': case L'“':
        case L'”': case L'‹': case L'›': case L'—':
        case L'–': case L'…':
            return 1;
        default:
            return 0;
    }
}
\end{lstlisting}

\vspace{2em}
\noindent\textbf{Нормализация текста.} Библиотека \texttt{normalizer\_lib.c} реализует функцию нормализации текста на широких символах wchar\_t. Нормализация выполняется путем преобразования текста в нижний регист и замены буквы "ё" на "е"

\begin{enumerate}
    
    \item \textbf{Проходим по тексту} символ за символом
    
    \item \textbf{Нормализуем символы:}
    \begin{itemize}
        \item Русские буквы в верхнем регистре преобразуются в нижний
        \item Английские буквы в верхнем регистре преобразуются в нижний
        \item Буква "ё"  преобразуется в "е"
        \item Цифры остаются без изменений
    \end{itemize}
\end{enumerate}

\begin{lstlisting}[language=C]
static
bool is_russian_letter( wchar_t c )
{
    if ( (c >= L'а' && c <= L'я') || 
         (c >= L'А' && c <= L'Я') || 
         (c == L'ё' || c == L'Ё') )
        return true;

    return false;
}

static
bool is_english_letter( wchar_t c )
{
    if ( (c >= L'a' && c <= L'z') || 
         (c >= L'A' && c <= L'Z') )
        return true;

    return false;
}

// normalization process in code
// =================================================
// ...
    if ( is_russian_letter( cur_char ) )
    {
        if (cur_char >= L'А' && cur_char <= L'Я')
            cur_char = L'а' + (cur_char - L'А');
        else if ( (cur_char == L'Ё') ||
                  (cur_char == L'ё') )
            cur_char = L'е';
    } else if ( is_english_letter( cur_char ) )
        cur_char = L'a' + (cur_char - L'A');
// ...
// =================================================
\end{lstlisting}

\vspace{2em}
\noindent\textbf{Удаление стоп-слов.} Библиотека \texttt{stopwords\_lib.c} реализует функцию удаления стоп-слов из текста на широких символах wchar\_t. Удаление выполняется путем выявления токенов, входящих в фиксированный список стоп-слов. Важно отметить, что тут и далее в процессе подготовки текста было принято решение оставлять английские токены без изменений. В нашем случае (статьи видеоигровой тематики на русском языке) англйиские слова редки, и как правило являются названиями проектов.

\begin{enumerate}
    
    \item \textbf{Проходим по тексту} токен за токеном
    
    \item \textbf{Фильтруем слова:}
    \begin{itemize}
        \item Если слово присутствует в списке стоп-слов – оно удаляется
        \item Если слово отсутствует в списке – оно сохраняется в выходном буфере
    \end{itemize}
\end{enumerate}

\begin{lstlisting}[language=C]
static const wchar_t* RUSSIAN_STOPWORDS[] =
{
    L"и", L"в", L"во", L"не", L"что", L"он", L"на", L"я", L"с", L"со",
    L"как", L"а", L"то", L"все", L"она", L"так", L"его", L"но", L"да",
    L"ты", L"к", L"у", L"же", L"вы", L"за", L"бы", L"по", L"только",
    L"ее", L"мне", L"было", L"вот", L"от", L"меня", L"еще", L" нет",
    L"о", L"из", L"ему", L"теперь", L"когда", L"даже", L"ну", L"ли",
    L"если", L"уже", L"или", L"ни", L"быть", L"был", L"него", L"до",
    L"вас", L"нибудь", L"опять", L"уж", L"вам", L"ведь", L"там", L"потом",
    L"себя", L"ничего", L"ей", L"может", L"они", L"тут", L"где", L"есть",
    L"надо", L"ней", L"для", L"мы", L"тебя", L"их", L"чем", L"была",
    L"сам", L"чтоб", L"без", L"будто", L"чего", L"раз", L"тоже", L"себе",
    L"под", L"будет", L"ж", L"тогда", L"кто", L"этот", L"того", L"потому",
    L"этого", L"какой", L"совсем", L"ним", L"здесь", L"этом", L"один",
    L"почти", L"мой", L"тем", L"чтобы", L"нее", L"сейчас", L"были",
    L"куда", L"зачем", L"всех", L"никогда", L"можно", L"при", L"наконец",
    L"два", L"об", L"другой", L"хоть", L"после", L"над", L"больше",
    L"тот", L"через", L"эти", L"нас", L"про", L"всего", L"них", L"какая",
    L"много", L"разве", L"три", L"эту", L"моя", L"впрочем", L"хорошо",
    L"свою", L"этой", L"перед", L"иногда", L"лучше", L"чуть", L"том",
    L"нельзя", L"такой", L"им", L"более", L"всегда", L"конечно", L"всю",
    L"между"
};

static const int RUSSIAN_STOPWORDS_COUNT = sizeof(RUSSIAN_STOPWORDS) / sizeof(RUSSIAN_STOPWORDS[0]);

static
bool is_stopword( const wchar_t *word )
{
    if ( !word || wcslen(word) == 0 )
        return false;
    
    for ( int i = 0; i < RUSSIAN_STOPWORDS_COUNT; i++ )
    {
        if ( wcscmp( word, RUSSIAN_STOPWORDS[i] ) == 0 )
            return true;
    }
    
    return false;
}
\end{lstlisting}

\vspace{2em}
\noindent\textbf{Результаты.} В результате токенизации было получено \textbf{22 364 175} токенов. Cредняя длина одного токена составила \textbf{6.49} символа.

Производительность:
\begin{itemize}
\item \textbf{Скорость обработки данных:} \textbf{2.24} Мб/сек.
\item \textbf{Скорость токенизации:} \textbf{135 996} токенов/сек.
\end{itemize}

Достигнутую производительность нельзя считать оптимальной. За счет внедрения параллельной обработки файлов можно было бы достигнуть кратного улучшения производительности.

Общее время работы полного пайплайна составило \textbf{164.45} секунды.

\vspace{2em}
\noindent\textbf{Результат токенизации.}
\begin{tcolorbox}[colback=blue!5!white,colframe=blue!75!black]
Where Winds Meet: Обзор самой эпической игры года Где встречаются The Witcher 3, Genshin Impact, Assassinss Creed и Ghost of Tsushima Игру «Там, где встречаются ветра» стоило бы назвать «Там, где встречаются тренды из всех игр сразу». В этом китайском чуде от Everstone Studio и NetEase Games можно разглядеть не только те проекты, которые я вынес в подзаголовок обзора, но и другие, включая Black Myth: Wukong и Sekiro: Shadows Die Twice. Но вот что удивительно — при этом в Where Winds Meet полно и своих фишек. Ну где ещё вы можете вступить во фракцию воинов-алкоголиков, сыграть в шахматы с лидером бандитского лагеря, кидаться медведями, драться с гусями, кататься на них верхом, а потом лечить гусей (нет, не от ПТСР, но было бы логично).
\end{tcolorbox}

\bigskip

\begin{tcolorbox}[colback=green!5!white,colframe=green!75!black]
where winds meet обзор самой эпической игры года встречаются the witcher 3 genshin impact assassins creed ghost of tsushima игру встречаются ветра стоило назвать встречаются тренды игр сразу китайском чуде everstone studio netease games разглядеть те проекты которые вынес подзаголовок обзора другие включая black myth wukong sekiro shadows die twice удивительно where winds meet полно своих фишек можете вступить фракцию воинов алкоголиков сыграть шахматы лидером бандитского лагеря кидаться медведями драться гусями кататься верхом лечить гусей нет птср логично
\end{tcolorbox}

\pagebreak

